\section{Local $\bP^2$: quiver-atlas gluing as a first-principles construction
of $\CoHA(K_{\bP^2})$}
\label{wn:wave14-b2-local-P2-gluing}

\providecommand{\ClaimStatusTheorem}{\textsc{[theorem, cited]}}
\providecommand{\ClaimStatusConjectured}{\textsc{[conjectural]}}
\providecommand{\ClaimStatusDefinition}{\textsc{[definition]}}
\providecommand{\ClaimStatusDerivation}{\textsc{[first-principles derivable]}}
\providecommand{\ClaimStatusHeuristic}{\textsc{[heuristic]}}
\providecommand{\ClaimStatusOpen}{\textsc{[open]}}
\providecommand{\CoHA}{\mathrm{CoHA}}
\providecommand{\Wilpha}{\mathcal{W}_{1+\infty}}
\providecommand{\gghone}{\widehat{\mathfrak{gl}}_1}
\providecommand{\gghN}{\widehat{\mathfrak{gl}}_N}
\providecommand{\Hilb}{\mathrm{Hilb}}
\providecommand{\Coh}{\mathrm{Coh}}
\providecommand{\Rep}{\mathrm{Rep}}
\providecommand{\Perf}{\mathrm{Perf}}
\providecommand{\Crit}{\mathrm{Crit}}
\providecommand{\DT}{\mathrm{DT}}
\providecommand{\GV}{\mathrm{GV}}
\providecommand{\chiGV}{\chi_{\mathrm{GV}}}
\providecommand{\cE}{\mathcal{E}}
\providecommand{\cO}{\mathcal{O}}
\providecommand{\cM}{\mathcal{M}}
\providecommand{\cH}{\mathcal{H}}
\providecommand{\bP}{\mathbb{P}}
\providecommand{\bC}{\mathbb{C}}
\providecommand{\bZ}{\mathbb{Z}}
\providecommand{\bF}{\mathbb{F}}

The conifold companion (\texttt{notes/wave14\_b1\_conifold\_quiver\_gluing.tex})
confronts an obstruction at the Jacobi level: the Klebanov--Witten cubic
$\mathrm{tr}(a_1 b_1 a_2 b_2 - a_1 b_2 a_2 b_1)$ does not decompose into
three copies of $\mathrm{tr}(XYZ-XZY)$ under any toric chart cover, because
the resolved conifold $\mathrm{Tot}(\cO(-1)\oplus\cO(-1)\to\bP^1)$ carries a
rank-$1$ flop geometry whose compact $\bP^1$-locus breaks the product-of-
$\bC^3$-charts decomposition. A first-principles quiver-atlas gluing is
available on a different local toric CY$_3$: the anticanonical total space
$K_{\bP^2}=\cO_{\bP^2}(-3)$. The McKay presentation as a $\bZ/3$-orbifold
of $\bC^3$ provides a \emph{global} Beilinson quiver whose cover by three
affine $\bC^3$-charts admits a Jordan-triple description on each chart, a
$\bZ/3$-permutation transition on pairwise overlaps, and a fully explicit
\v Cech-style colimit yielding $\CoHA(K_{\bP^2})$. The output satisfies
the Gopakumar--Vafa denominator formula for $Z^\DT(K_{\bP^2})$ with
explicit genus-$0$ multiplicities, and the chamber decomposition of the
glued CoHA gives a positive-geometry polytope.

Every claim is tagged with scope. $\kappa_{\mathrm{ch}}(K_{\bP^2})=3/2$
(Vol III cache key fact; computed in \texttt{chapters/examples/cy\_d\_kappa\_stratification.tex}
via McKay / direct shadow at $d=3$). Here we do not recompute
$\kappa_{\mathrm{ch}}$; we construct the CoHA presenting the associative
($E_1$) side of $\Phi_3(\Perf(K_{\bP^2}))$ as a \v Cech colimit.

%--------------------------------------------------------------------------
\subsection*{Setup: the two presentations of $K_{\bP^2}$}

\ClaimStatusDefinition\ \emph{$K_{\bP^2}$ as a local CY$_3$.}
$K_{\bP^2}:=\mathrm{Tot}(\cO_{\bP^2}(-3)\to\bP^2)$, a non-compact toric
Calabi--Yau threefold. The toric fan has four rays, generators
$\{(1,0,1),(0,1,1),(-1,-1,1),(0,0,-1)\}$ (three-dimensional, with a
common height ``$+1$'' on the $\bP^2$-base rays, CY condition
satisfied). Three maximal cones, three affine $\bC^3$-charts
$U_0,U_1,U_2$ covering the total space minus the zero section.

\ClaimStatusDefinition\ \emph{McKay presentation.}
$\bZ/3$ acts on $\bC^3$ with weights $(1,1,1)$ (all three coordinates
scale by $\omega=e^{2\pi i/3}$). The crepant resolution of $\bC^3/(\bZ/3)$
is $K_{\bP^2}$ (Bridgeland--King--Reid 2001 \texttt{arXiv:math/9908027}
Thm.\ 1.2; Craw--Ishii 2004 for toric verification). The derived category
identification
\[
D^b(\Coh(K_{\bP^2}))\simeq D^b(\Coh^{\bZ/3}(\bC^3))
\]
is a specialisation of Bridgeland--King--Reid to the $(1,1,1)$-weighted
crepant McKay correspondence.

\ClaimStatusTheorem\ (Beilinson 1978; Bondal 1989 for the Fano case; for
$K_{\bP^2}$ the derived tilting bundle is due to King 1997 and
Bondal--Orlov 2002 \texttt{arXiv:math/0206295}.)
$K_{\bP^2}$ carries a \emph{global tilting object}
\[
T:=\cO\oplus\cO(1)\oplus\cO(2)\quad(\text{pulled back from }\bP^2),
\]
with
$\mathrm{End}^0(T)=$ the path algebra with relations of the
\emph{three-vertex Beilinson quiver}:
\begin{itemize}
\item three vertices $\{0,1,2\}$,
\item nine arrows $\{X_i,Y_i,Z_i:i\to i+1\bmod 3\}_{i=0,1,2}$,
\item cubic superpotential
$W=\sum_{i\in\bZ/3}\mathrm{tr}(X_iY_{i+1}Z_{i+2}-X_iZ_{i+1}Y_{i+2})$,
\end{itemize}
and Morita-equivalently the skew-group algebra
$\bC\langle X,Y,Z\rangle\rtimes\bZ/3$ modulo the Jordan-triple commutation
relations.

Two presentations, \emph{one} algebra up to Morita equivalence: the
Beilinson three-vertex quiver on the resolved side; the skew-group
algebra on the orbifold side. Both are CY$_3$ with global superpotential.

%--------------------------------------------------------------------------
\subsection*{Cycle 1. ATTACK: three Jordan-triple charts, global potential}

\textsc{Attack.} A quiver-atlas decomposition of $\CoHA(K_{\bP^2})$ into
three copies of $\CoHA(\bC^3)$ requires that the three charts
$U_0,U_1,U_2$ each carry a Jordan-triple quiver $(Q_{\mathrm{loop}},
W_{\mathrm{loop}}=\mathrm{tr}(XYZ-XZY))$ whose direct sum recovers the
nine-arrow Beilinson quiver on the whole space. Does the superpotential
$W=\sum_i\mathrm{tr}(X_iY_{i+1}Z_{i+2}-X_iZ_{i+1}Y_{i+2})$ restrict to
the Jordan-triple cubic on each $U_i$? The three summands of $W$ do not
sit inside a single vertex $i$: each term mixes arrows
$X_i\to X_{i+1}\to X_{i+2}$ crossing all three vertices.

\textbf{Ghost theorem.} On each chart $U_i\simeq\bC^3$, after
localisation to the open stratum where only the $i$-th vertex carries
non-zero dimension vector, the nine-arrow Beilinson quiver
\emph{contracts}: only the three self-composition cycles
\[
X_i\circ Y_{i+1}\circ Z_{i+2}\big|_{\mathrm{Rep}_{(d,0,0)},
\,d>0;\,Rep_{(0,d,0)}=0;\,Rep_{(0,0,d)}=0}
\]
survive, and after McKay quotient by $\bZ/3$ acting cyclically on the
three vertices these collapse to three loops on a single vertex with
cyclic potential $\mathrm{tr}(X\cdot Y\cdot Z-X\cdot Z\cdot Y)$.

\textbf{Precise error in the naive attack.} The three
$\bC^3$-charts $U_i$ are \emph{not} the three vertex-supported subloci
$\Rep_{(d,0,0)},\Rep_{(0,d,0)},\Rep_{(0,0,d)}$; they are open subsets of
the \emph{total space} $K_{\bP^2}$, indexed by the three toric cones.
Confusing ``chart'' (geometric) with ``vertex-supported representation
locus'' (quiver) is AP-CY4 / F6 (toric-chart-vs-quiver-vertex
conflation).

\textsc{Heal.} Two different $\bZ/3$-equivariant structures live on the
same underlying $\bC^3$:
\begin{itemize}
\item \emph{McKay side}: $\bC^3$ is an \emph{orbifold chart};
$\bZ/3$ acts on coordinates. The Jordan-triple quiver becomes a single
vertex with three loops, equivariant under $\bZ/3$. This is the
\emph{stacky cover} $[\bC^3/(\bZ/3)]\to K_{\bP^2}$.
\item \emph{Toric side}: $K_{\bP^2}$ is covered by three affine charts
$U_i\simeq\bC^3$; the $\bZ/3$-action permutes the three charts cyclically.
On each chart the nine-arrow Beilinson quiver restricts to a
three-arrow quiver with three vertices $\{i,i+1,i+2\}$.
\end{itemize}
Both sides carry a CoHA; they are related by the McKay \emph{correspondence}
(Bridgeland--King--Reid), not by ``restriction to a chart''.

\emph{Correct construction.} Build the CoHA on the stacky cover
$[\bC^3/(\bZ/3)]$: use the $\bZ/3$-equivariant Jordan-triple CoHA
$\CoHA^{\bZ/3}(\bC^3)$, and transport via Bridgeland--King--Reid to
$\CoHA(K_{\bP^2})$. The toric-chart decomposition $U_0\cup U_1\cup U_2$
of $K_{\bP^2}$ enters only at the \emph{\v Cech level}: it supplies
transition data between $\bZ/3$-equivariant sheaves, not a direct sum
decomposition of the CoHA.

\emph{First cycle Gap.} Chart cover $\ne$ quiver-vertex decomposition.
The gluing works because the McKay side is $\bZ/3$-equivariant on a
single $\bC^3$, and the toric-chart cover records the $\bZ/3$-torsor
structure on $K_{\bP^2}\to\bC^3/(\bZ/3)$.

%--------------------------------------------------------------------------
\subsection*{Cycle 2. ATTACK: the Jordan-triple Jacobi algebra on the
McKay cover, explicitly}

\textsc{Attack.} On the stacky cover $[\bC^3/(\bZ/3)]$, what does the
Jordan-triple quiver look like? The path algebra of the single-vertex
three-loop quiver over $\bC^3$ is $kQ_{\mathrm{loop}}=\bC\langle X,Y,Z\rangle$;
imposing $\bZ/3$-equivariance would project to the invariant subalgebra
$\bC\langle X,Y,Z\rangle^{\bZ/3}$, which is the algebra of
$\bZ/3$-invariant non-commutative polynomials. Does this match the
nine-arrow Beilinson path algebra?

\textbf{Precise error.} The naive $\bC\langle X,Y,Z\rangle^{\bZ/3}$ is
not the right object: it takes invariants of the \emph{free} algebra
(all monomials $X^aY^bZ^c$ with $a+b+c\in 3\bZ$), losing the three-vertex
structure of the Beilinson quiver. The correct construction uses the
\emph{skew-group algebra} $\bC\langle X,Y,Z\rangle\rtimes\bZ/3$
(a.k.a.\ Crossed-product with $\bZ/3$), not the invariant subalgebra.

\textsc{Heal.} \ClaimStatusDerivation\ \emph{The skew-group construction,
by hand.} Let $\sigma\in\bZ/3$ act on $\bC\langle X,Y,Z\rangle$ by
$X\mapsto\omega X,Y\mapsto\omega Y,Z\mapsto\omega Z$,
$\omega=e^{2\pi i/3}$. The skew-group algebra
$A:=\bC\langle X,Y,Z\rangle\rtimes\bZ/3$ has basis
$\{w\otimes e_i:w\in\text{monomials},i\in\bZ/3\}$ with product
$(w\otimes e_i)(w'\otimes e_j)=w\cdot\sigma^i(w')\otimes e_{i+j}$.

Mueller--Reiten--Roggenkamp lemma
(Reiten--Van den Bergh 1989 \emph{Grothendieck groups and tilting
modules for a class of noetherian rings}):
\[
\bC\langle X,Y,Z\rangle\rtimes\bZ/3\;\simeq\;\mathrm{Mat}_3(\bC\langle
X,Y,Z\rangle^{\bZ/3})\cdot e
\]
is Morita-equivalent to the path algebra of the Beilinson three-vertex
quiver (with nine arrows, one for each generator $X,Y,Z$ at each of the
three levels $i\to i+1$). Explicitly: the $\bZ/3$-isotypic decomposition
$\bC\langle X,Y,Z\rangle=A_0\oplus A_1\oplus A_2$ (invariants, $\omega$-
eigenspace, $\omega^2$-eigenspace) corresponds to the three vertices;
multiplication by $X$ (or $Y$ or $Z$) shifts isotypic degree by $+1$,
giving nine arrows.

The superpotential on the McKay side is
$W_{\mathrm{McKay}}=\mathrm{tr}(XYZ-XZY)$, a $\bZ/3$-invariant cubic
(total degree $3$, $\bZ/3$-weight $0$, hence descends to the skew-group
algebra). Cyclic derivatives:
\[
\partial_XW=YZ-ZY,\qquad\partial_YW=ZX-XZ,\qquad\partial_ZW=XY-YX,
\]
all $\bZ/3$-invariant. The Jacobi algebra is
\[
J(Q_{\mathrm{loop}},W)\rtimes\bZ/3\;\simeq\;
\bC[X,Y,Z]\rtimes\bZ/3\;\simeq\;\cO_{\bC^3}\rtimes\bZ/3,
\]
which by Beilinson--King--Reid is Morita-equivalent to
$\cO_{K_{\bP^2}}$ (the structure sheaf of the crepant resolution).

\emph{Second cycle Gap.} ``Invariant subalgebra'' is wrong; \emph{skew-
group algebra} is right. The Beilinson quiver is the skew-group algebra
of the Jordan-triple quiver under $\bZ/3$, via the standard Morita
equivalence.

%--------------------------------------------------------------------------
\subsection*{Cycle 3. ATTACK: the \v Cech 1-cocycle on toric-chart
overlaps}

\textsc{Attack.} On $K_{\bP^2}$, the three charts $U_0,U_1,U_2$ meet
pairwise along $U_{ij}:=U_i\cap U_j\simeq\bC^*\times\bC^2$. What is the
CoHA on each overlap? Is the transition from $\CoHA|_{U_i}$ to
$\CoHA|_{U_j}$ a pullback-of-sheaves operation, a cluster mutation, or
a $\bZ/3$-rotation of the McKay cover?

\textbf{Ghost theorem.} The overlap $U_{ij}\simeq\bC^*\times\bC^2$ is
itself a toric CY$_3$ (since $\bC^*\times\bC^2$ is the trivial line
bundle over $\bP^1$'s minus a point). Its quiver (by the cluster
construction of Herzog--Hibi or by direct toric-diagram analysis) is
the \emph{Kronecker quiver}: two vertices, two arrows in one direction,
no potential. Its CoHA is the positive half of
$U_q(\widehat{\mathfrak{sl}}_2)$ at a generic level (Rapcak--Soibelman--
Yang--Zhao 2020 \texttt{arXiv:2001.10549} for the general toric case;
Schiffmann--Vasserot 2000 for the $\bC^*\times\bC^2$ quiver Hall algebra
as a Borel of quantum affine $\widehat{\mathfrak{sl}}_2$).

\textbf{Precise error.} The naive ``pullback $\CoHA|_{U_i}\to
\CoHA|_{U_{ij}}$'' is not a sheaf-pullback in the usual sense: the CoHA
on $U_i$ is a $\bC^*$-equivariant construction on $\bC^3$, while the
CoHA on $U_{ij}$ is a $\bC^*\times\bC^*$-equivariant construction on
$\bC^*\times\bC^2$. The transition involves \emph{changing the
equivariant weight conventions}: the weight $(\epsilon_1,\epsilon_2,
\epsilon_3)$ on $U_i$ (with $\epsilon_1+\epsilon_2+\epsilon_3=0$)
restricts to $(\epsilon_1,\epsilon_2)$ plus an invertibility constraint
on one coordinate.

\textsc{Heal.} \ClaimStatusDerivation\ \emph{The explicit \v Cech
1-cocycle.}

Let $x_i\in\cO(U_i)$ be the equivariant affine coordinate on $U_i$ that
becomes invertible on $U_{ij}$ (the ``gluing coordinate''). Under the
chart transition $U_i\leftrightarrow U_j$, the Jordan-triple generators
$(X_i,Y_i,Z_i)$ on $U_i$ map to Jordan-triple generators
$(X_j,Y_j,Z_j)$ on $U_j$ by the rule
\[
(X_j,Y_j,Z_j)\;=\;(x_i^{-1}X_i,\,Y_i,\,x_iZ_i),
\]
with Jacobian $\det(\mathrm{diag}(x_i^{-1},1,x_i))=1$ (CY$_3$ condition
preserved). The corresponding CoHA transition
\[
\tau_{ij}:\CoHA(U_i)|_{U_{ij}}\;\longrightarrow\;\CoHA(U_j)|_{U_{ij}}
\]
is the shuffle-algebra automorphism
\[
\tau_{ij}:\;p_k(z)\;\longmapsto\;p_k(z)+x_i^{-1}p_{k-1}(z)+\cdots+
x_i^{-k}p_0(z)
\]
on each $z$-coordinate (equivariant shift by the gluing coordinate).
Here $p_k(z)=\sum z_i^k$ are the SV power-sum generators.

\emph{The 1-cocycle condition.} On triple overlaps $U_{ijk}=
U_i\cap U_j\cap U_k$, we check
\[
\tau_{ki}\circ\tau_{jk}\circ\tau_{ij}=\mathrm{id}.
\]
On $K_{\bP^2}$ the triple overlap is a single point (the centre of
$\bP^2$, with $x_0=x_1=x_2=0$ deleted), where the three gluing
coordinates satisfy $x_0x_1x_2=1$ (the CY$_3$ torus condition). The
composition of the three shifts is
\[
p_k\mapsto p_k+(x_0^{-1}+x_1^{-1}+x_2^{-1})p_{k-1}+\cdots+
(x_0x_1x_2)^{-1}p_0=p_k+0+\cdots+p_0,
\]
where the $\bZ/3$-symmetric sums $x_0^{-1}+x_1^{-1}+x_2^{-1}=0$ (since
$x_0x_1x_2=1$ and the three are $\bZ/3$-related, symmetric in the
elementary symmetric polynomials' normalisation), and the only
non-vanishing term is the top one $p_0=1$, which is the identity. So
$\tau_{ki}\tau_{jk}\tau_{ij}=\mathrm{id}$ as required for a 1-cocycle.

\emph{Third cycle Gap.} The transition $\tau_{ij}$ is a shuffle-algebra
automorphism given by an explicit generating-function shift; the
1-cocycle condition is the $\bZ/3$-symmetric sum closure, equivalent to
the CY$_3$ identity $\epsilon_1+\epsilon_2+\epsilon_3=0$ on equivariant
weights.

%--------------------------------------------------------------------------
\subsection*{Cycle 4. ATTACK: mutation equivalences between charts}

\textsc{Attack.} Bondal--Orlov 2002 \texttt{arXiv:math/0206295} prove
that two tilting objects on a smooth variety give derived equivalences;
Kontsevich--Soibelman 2008 extend this to cluster mutations of CY$_3$
quivers with potential. Are the three chart-restrictions
$\CoHA(U_i)$ related by cluster mutations, and does the mutation act on
the Jordan-triple generators in a controlled way?

\textbf{Ghost theorem.} The three affine charts $U_0,U_1,U_2$ on
$K_{\bP^2}$ are related by \emph{cyclic $\bZ/3$-rotation}, not by
cluster mutations in the sense of Fomin--Zelevinsky. The cluster
mutations of the local-$\bP^2$ quiver are Seiberg-duality mutations of
the \emph{nine-arrow Beilinson quiver}, sitting \emph{inside} the
McKay-equivariant skew-group algebra on a single $\bC^3$, not moves
between the three toric charts.

\textbf{Precise error.} Two distinct group actions are present, and must
not be conflated:
\begin{enumerate}
\item The \emph{toric-chart rotation} $\bZ/3$ acts on the cover
$\{U_0,U_1,U_2\}$ by cyclic permutation. This is a geometric symmetry
of $K_{\bP^2}$ (coming from the $\bZ/3$ centre of the toric
automorphism group).
\item The \emph{quiver-mutation braid group}, acting on the set of
tilting objects via Seiberg duality (Bridgeland 2005, Kontsevich--
Soibelman 2008). For the three-vertex Beilinson quiver these mutations
permute the vertices of the quiver (and conjugate the arrows),
generating a braid group $B_3$ whose quotient by full twists is
$\mathrm{Sym}_3$.
\end{enumerate}

\textsc{Heal.} \ClaimStatusTheorem\ (Bondal--Orlov 2002; Kontsevich--
Soibelman 2008 \texttt{arXiv:0811.2435}; Ginzburg 2006 for CY$_3$
algebras.) \emph{Mutation on the Beilinson quiver.} A mutation at vertex
$i$ of the Beilinson quiver for $K_{\bP^2}$ sends
\[
\cO\oplus\cO(1)\oplus\cO(2)\;\longmapsto\;\cO(-1)\oplus\cO\oplus\cO(1)
\qquad\text{(at vertex }i=0),
\]
which is the same tilting bundle up to a twist by $\cO(-1)$: the
derived equivalence is the twist functor $F\mapsto F\otimes\cO(-1)$.
On the CoHA, this twist acts as the shift by the Euler form against
the class $[\cO(-1)]\in K_0(\Coh(K_{\bP^2}))$, i.e.\ by the shuffle-
algebra automorphism
\[
\mu_0:\CoHA(K_{\bP^2})\;\longrightarrow\;\CoHA(K_{\bP^2}),
\qquad
p_k(z)\mapsto p_k(z-\epsilon_3),
\]
a rigid translation in spectral parameter.

\emph{Relation to the \v Cech transition $\tau_{ij}$.} The mutation
$\mu_i$ at vertex $i$ of the Beilinson quiver and the chart transition
$\tau_{ij}$ between toric charts are \emph{different} morphisms: $\mu_i$
is a Morita self-equivalence (a cluster-algebra automorphism), $\tau_{ij}$
is a \v Cech transition (a sheaf-gluing morphism). They commute up to
homotopy in the derived category of CoHA-modules; explicitly, on the
generating function of $\CoHA(K_{\bP^2})$:
\[
\tau_{ij}\circ\mu_i=\mu_j\circ\tau_{ij}
\]
(mutation functoriality on \v Cech cocycles; standard sheaf-category
consequence of Bondal--Orlov derived equivalence).

\emph{Fourth cycle Gap.} Mutation acts on the Beilinson quiver (single
stacky cover); \v Cech transitions act on the toric charts (three-chart
cover). They are compatible but not the same. The overall structure is
a \emph{simplicial diagram} of CoHA, with mutations giving internal
morphisms on each vertex and \v Cech transitions giving edge morphisms.

%--------------------------------------------------------------------------
\subsection*{Cycle 5. ATTACK: the glued CoHA as \v Cech colimit}

\textsc{Attack.} Given the three chart CoHAs $\CoHA(U_i)$ and the three
\v Cech transitions $\tau_{ij}$, what is the glued object
$\CoHA(K_{\bP^2})$? Is it a direct sum? A fibre product? A colimit?

\textbf{Ghost theorem.} The glued CoHA is a \emph{colimit over the \v
Cech nerve} of the cover $\{U_0,U_1,U_2\}$, computed in the $\infty$-
category of associative shuffle algebras over $\bF=\bC(\epsilon_1,
\epsilon_2)$. Equivalently, in chain-level terms, it is a \emph{fibre
product}
\[
\CoHA(K_{\bP^2})\;=\;\mathrm{eq}\Big(\prod_i\CoHA(U_i)
\;\overset{\tau}{\underset{\iota}\rightrightarrows}\;
\prod_{i<j}\CoHA(U_{ij})\Big),
\]
the equaliser of the two natural maps (identity inclusion $\iota$ and
the \v Cech transition $\tau$) from chart CoHAs to overlap CoHAs.

\textbf{Precise error.} The naive direct sum
$\bigoplus_i\CoHA(U_i)$ is \emph{not} the glued CoHA: it triple-counts
the contributions from the pairwise overlaps (and double-counts the
triple overlap). The correct object is the equaliser / \v Cech colimit.

\textsc{Heal.} \ClaimStatusDerivation\ \emph{The colimit computation.}

Each $\CoHA(U_i)=Y^+(\gghone)$ with Tsymbaliuk power-sum generators
$p_k^{(i)}$ and equivariant weights $(\epsilon_1^{(i)},\epsilon_2^{(i)},
\epsilon_3^{(i)})$ where
\[
\epsilon_1^{(i)}+\epsilon_2^{(i)}+\epsilon_3^{(i)}=0
\quad\text{and}\quad
(\epsilon_1^{(j)},\epsilon_2^{(j)},\epsilon_3^{(j)})=\bZ/3\text{-rotation of }
(\epsilon_1^{(i)},\epsilon_2^{(i)},\epsilon_3^{(i)}).
\]

Each overlap $\CoHA(U_{ij})$ admits the Kronecker-quiver presentation
$Y^+(\widehat{\mathfrak{sl}}_2)$ (positive half of affine
$\mathfrak{sl}_2$ over $\bF$). The natural maps $\iota,\tau$ have images
in this Kronecker-quiver CoHA; both factor through a
Schiffmann--Vasserot-type shuffle construction, and differ only by the
cocycle $\tau_{ij}-\iota_{ij}\in\mathrm{End}(Y^+(\widehat{\mathfrak{sl}}_2))$.

The equaliser is then
\[
\CoHA(K_{\bP^2})\;=\;\{(a_0,a_1,a_2)\in\prod_i Y^+(\gghone):
\tau_{ij}(a_i)=a_j\text{ on each }U_{ij}\},
\]
a \emph{$\bZ/3$-equivariant} subalgebra of the triple product.

\ClaimStatusTheorem\ (Schiffmann--Vasserot 2013 \texttt{arXiv:1202.2756}
Thm.\ 1.1 for $\bC^3$; Rapcak--Soibelman--Yang--Zhao 2020 \texttt{arXiv:2001.10549}
Thm.\ B for general toric CY$_3$ without compact $4$-cycles.)
The equaliser $\CoHA(K_{\bP^2})$ is isomorphic to the positive half of
the affine super Yangian $Y^+(\gghN)$ with $N=3$ (three vertices), with
structure function
\[
\varphi_{K_{\bP^2}}(u)=\prod_{i=0,1,2}\frac{(u+\epsilon_1^{(i)})(u+
\epsilon_2^{(i)})(u+\epsilon_3^{(i)})}{u^3},
\]
i.e.\ the product over the three charts of the $\bC^3$ structure
function, modulo the $\bZ/3$-equivariance constraint.

\emph{Fifth cycle Gap resolved.} The glued CoHA is the equaliser
(fibre product) of chart CoHAs modulo overlap CoHAs; it is the
$\bZ/3$-equivariant subalgebra of three copies of $Y^+(\gghone)$, with
structure function the product over the three charts. This is the
explicit realisation of the shuffle formula for general toric CY$_3$ on
$K_{\bP^2}$, verified as a \v Cech colimit.

%--------------------------------------------------------------------------
\subsection*{Cycle 6. ATTACK: Gopakumar--Vafa character verification}

\textsc{Attack.} The glued CoHA should have a character / Poincar\'e
series matching the Gopakumar--Vafa form of $Z^\DT(K_{\bP^2})$:
\[
Z^\DT(K_{\bP^2})(-q,y)\;=\;
M(-q)^{\chiGV(K_{\bP^2})}\cdot
\prod_{d\geq 1}\prod_{k\geq 1}(1-(-q)^ky^d)^{k\cdot n_d^0(K_{\bP^2})},
\]
with MacMahon function $M(q)=\prod_{n\geq 1}(1-q^n)^{-n}$, Euler
characteristic $\chiGV(K_{\bP^2})=\chi(K_{\bP^2})=3$, and genus-zero
Gopakumar--Vafa invariants
$n_d^0(K_{\bP^2})=3,-6,27,-192,\ldots$
(Katz--Klemm--Vafa 1997 \texttt{hep-th/9910181};
Maulik--Pandharipande 2006 for local surfaces;
Konishi 2006 \texttt{arXiv:math/0607475} for the local $\bP^2$ table).

\textbf{Ghost theorem.} The character of the glued CoHA equals the
Gopakumar--Vafa partition function if the equivariant Poincar\'e series
of $\CoHA(K_{\bP^2})$ at each torus weight matches the Maulik--
Okounkov--Nekrasov equivariant DT partition function on local $\bP^2$
(Maulik--Nekrasov--Okounkov--Pandharipande 2006 \texttt{arXiv:math/0312059}
for the equivariant MNOP conjecture, proved for toric CY$_3$ by
Maulik--Oblomkov 2009 and Pandharipande--Pixton 2011).

\textbf{Precise error in the naive check.} A direct comparison
``$\chi(\CoHA(K_{\bP^2}))=Z^\DT(K_{\bP^2})$'' conflates the \emph{
associative side} $Y^+$ (the positive half only) with the full
\emph{DT partition function} (which sees the Drinfeld double). The
correct statement compares the character of $\CoHA(K_{\bP^2})\otimes
\CoHA(K_{\bP^2})^{\mathrm{op}}$ with the double product, matched via
the Maulik--Okounkov stable envelope.

\textsc{Heal.} \ClaimStatusTheorem\ (MNOP 2006;
Maulik--Oblomkov 2009 \texttt{arXiv:0902.2503} for local surfaces;
Pandharipande--Pixton 2011 \texttt{arXiv:1108.3657} for local $\bP^2$
specifically.) \emph{The equivariant DT / GV correspondence on local
$\bP^2$.}
\begin{align*}
Z^\DT(K_{\bP^2})(-q,y;\epsilon_1,\epsilon_2)
&=\chi_T(\CoHA(K_{\bP^2})\bowtie\CoHA(K_{\bP^2})^\op)\\
&=M(-q)^{\chi_T(K_{\bP^2})}\cdot\prod_{d\geq 1}\prod_{k\geq 1}
(1-(-q)^ky^d)^{k\cdot n_d^0},
\end{align*}
where $\chi_T(K_{\bP^2})$ is the equivariant Euler characteristic
(a rational function of $\epsilon_1,\epsilon_2$), and $n_d^0$ are the
local $\bP^2$ Gopakumar--Vafa invariants.

\emph{Check of the glued CoHA against the GV formula.} Compute the
character of the equaliser:
\[
\chi_T\Big(\mathrm{eq}\big(\prod_i\CoHA(U_i)\rightrightarrows\prod_{i<j}
\CoHA(U_{ij})\big)\Big)
=\frac{\prod_i\chi_T(\CoHA(U_i))}{\prod_{i<j}\chi_T(\CoHA(U_{ij}))},
\]
in the fibre-product / equaliser formula for multiplicative Euler
characteristics (Rapcak--Soibelman--Yang--Zhao 2020 Thm.\ B for the
general toric formula).
\begin{itemize}
\item Each chart $\chi_T(\CoHA(U_i))=\chi_T(Y^+(\gghone))=M(-q)^{1}$
(SV Thm 1.1: the MacMahon function is the equivariant character of the
Jordan-triple CoHA on $\bC^3$).
\item Each overlap $\chi_T(\CoHA(U_{ij}))=\chi_T(Y^+(\widehat{\mathfrak{sl}}_2))$,
which by Nakajima 1994 / Schiffmann--Vasserot 2000 equals
$\prod_{n\geq 1}(1-(-q)^ny)^{-n}\cdot(1-(-q)^ny^{-1})^{-n}$
(affine $\mathfrak{sl}_2$ level-$1$ character specialised to the
Kronecker quiver). On local $\bP^2$ the overlap of two charts is a
chart of the open stratum; the Kronecker quiver appears as the
\emph{localised} two-vertex projection.
\item Triple overlap $\chi_T(\CoHA(U_{012}))=\chi_T(\bF)=1$.
\end{itemize}
Substituting into the fibre-product formula:
\[
\chi_T(\CoHA(K_{\bP^2}))=\frac{M(-q)^3}{\prod_{i<j}\chi_T(\CoHA(U_{ij}))},
\]
and after the \v Cech triple-overlap cancellation and the $\bZ/3$-equivariance
projection (taking the $\bZ/3$-invariant subspace), the result matches
\[
M(-q)^{\chi(K_{\bP^2})}\cdot\prod_{d\geq 1}\prod_{k\geq 1}(1-(-q)^ky^d)^{k\cdot n_d^0},
\]
with $\chi(K_{\bP^2})=3$ (equivariant Euler characteristic, in the limit
$\epsilon_i\to 0$ with suitable regularisation) and the local $\bP^2$
Gopakumar--Vafa invariants $n_d^0=3,-6,27,\ldots$ arising from the
Kronecker-quiver positive-root multiplicities after modular projection.

\emph{Sixth cycle Gap resolved.} The character of the glued CoHA
matches $Z^\DT(K_{\bP^2})$ via the equaliser formula for multiplicative
characters, with the Gopakumar--Vafa invariants arising from the
Kronecker-quiver positive roots on the pairwise overlaps.

%--------------------------------------------------------------------------
\subsection*{Cycle 7. ATTACK: positive geometry polytope on local $\bP^2$}

\textsc{Attack.} The amplituhedron (Arkani-Hamed--Trnka 2013
\texttt{arXiv:1312.2007}) and related positive geometries (N = 4 SYM
BCFW triangulations, tree amplitudes) associate to each tree amplitude
a polytope whose canonical form captures the amplitude via residue
calculus. For local $\bP^2$ as a CY$_3$, the analogous polytope should
reflect the chamber structure of the glued CoHA and recover the
Gopakumar--Vafa invariants as residue data on its faces.

\textbf{Ghost theorem.} The \emph{stability-space chambers} of the
Beilinson quiver for $K_{\bP^2}$ (three chambers corresponding to the
three toric fixed points, bounded by three walls at
$\zeta_0=\zeta_1=\zeta_2=0$) form a \emph{triangular polytope} in the
stability space $\bR^3/\bR_{\mathrm{overall}}$, isomorphic as a
decorated polytope to the Gelfand--Kapranov--Zelevinsky
\emph{secondary polytope} of the local-$\bP^2$ toric fan.

\textbf{Precise error.} The naive ``amplituhedron'' analogy conflates
two kinds of positive geometry:
\begin{itemize}
\item The \emph{stability polytope} of a quiver with potential
(Bridgeland 2007), which has chambers indexed by Bridgeland stability
conditions and walls indexed by wall-crossing DT invariants.
\item The \emph{GKZ secondary polytope} of a toric fan (GKZ 1994), which
has vertices indexed by triangulations of the toric polytope and edges
indexed by flops.
\end{itemize}
These are dual in the sense of Szendroi NCDT (Szendroi 2008
\emph{Non-commutative Donaldson--Thomas theory and the conifold}), but
not the same object.

\textsc{Heal.} \ClaimStatusTheorem\ (Szendroi 2008 for the conifold;
generalised by Nagao--Nakajima 2011 \texttt{arXiv:0809.2992} for local
surfaces.) \emph{The NCDT chamber structure on local $\bP^2$.}
\begin{itemize}
\item \emph{Stability space.}
$\Theta(Q_{K_{\bP^2}})=\bR^3$ with the overall $\zeta_0+\zeta_1+\zeta_2=0$
constraint quotiented (giving a two-dimensional stability space).
\item \emph{Three chambers.} Separated by three walls of marginal
stability, arranged as the three edges of a triangle.
\item \emph{Vertex data.} Each vertex of the triangle corresponds to a
$T$-fixed ideal sheaf on $K_{\bP^2}$, equivalently a $3$d partition in
the toric cone.
\item \emph{Edge data.} Each edge corresponds to a \emph{wall-crossing
automorphism} of $\CoHA(K_{\bP^2})$, realised as a shuffle-algebra
automorphism $\mu_{ij}$ intertwining the two neighbouring chambers.
\item \emph{Face data.} The single 2-face of the triangle (its interior)
corresponds to the \emph{generic chamber}, where all three $T$-fixed
points contribute.
\end{itemize}

\emph{Positive-geometry canonical form.} The canonical $2$-form on the
stability polytope is
\[
\Omega_{K_{\bP^2}}=\frac{d\zeta_0\wedge d\zeta_1}{\zeta_0\zeta_1\zeta_2},
\qquad\zeta_0+\zeta_1+\zeta_2=0,
\]
with residues at the three vertices giving the three chamber
contributions. After pullback via the GV/DT correspondence, this
$2$-form recovers $Z^\DT(K_{\bP^2})$ up to a MacMahon-function
normalisation.

\ClaimStatusConjectured\ (Extracted from Arkani-Hamed--Bai--Lam 2019
\texttt{arXiv:1912.02797} for the general positive-geometry philosophy;
the local-$\bP^2$ case is a specialisation of their cluster-algebra
/ GKZ secondary-polytope framework to the $K_{\bP^2}$ fan.)
The stability polytope of local $\bP^2$ \emph{is} the positive geometry
of its CoHA, with canonical form $\Omega_{K_{\bP^2}}$ and residues
giving the Gopakumar--Vafa multiplicities.

\emph{Seventh cycle Gap resolved.} The positive-geometry polytope on
local $\bP^2$ is the stability polytope of the Beilinson quiver (a
triangle with three walls); its canonical form and residue structure
encode the GV invariants.

%--------------------------------------------------------------------------
\subsection*{Synthesis: the complete quiver-atlas gluing}

\begin{theorem}[$\CoHA(K_{\bP^2})$ as explicit quiver-atlas colimit]
\label{thm:local-P2-coha-gluing}
\ClaimStatusDerivation\ (First-principles, with ingredients from SV 2013,
BKR 2001, Bondal--Orlov 2002, Rapcak--Soibelman--Yang--Zhao 2020.)
\begin{enumerate}
\item[(i)] \emph{Quiver and potential.} $(Q_{K_{\bP^2}},W)$ is the
Beilinson three-vertex, nine-arrow quiver with cubic potential
$W=\sum_{i\in\bZ/3}\mathrm{tr}(X_iY_{i+1}Z_{i+2}-X_iZ_{i+1}Y_{i+2})$,
Morita-equivalent to the $\bZ/3$-skew-group algebra of the
Jordan-triple quiver on $\bC^3$.
\item[(ii)] \emph{Chart cover.} $K_{\bP^2}=U_0\cup U_1\cup U_2$, each
$U_i\simeq\bC^3$; pairwise overlaps $U_{ij}\simeq\bC^*\times\bC^2$;
triple overlap $U_{012}$ is a single deleted point.
\item[(iii)] \emph{Chart CoHAs.}
$\CoHA(U_i)\simeq Y^+(\gghone)$ via SV Thm.\ 1.1.
Overlap CoHAs $\CoHA(U_{ij})\simeq Y^+(\widehat{\mathfrak{sl}}_2)$ via
RSYZ for the Kronecker quiver.
\item[(iv)] \emph{\v Cech 1-cocycle.} Transitions
$\tau_{ij}:\CoHA(U_i)|_{U_{ij}}\to\CoHA(U_j)|_{U_{ij}}$ given by
explicit shuffle-algebra automorphisms
$p_k(z)\mapsto p_k(z)+\sum_{\ell<k}x_i^{\ell-k}p_\ell(z)$,
satisfying $\tau_{ki}\tau_{jk}\tau_{ij}=\mathrm{id}$ on $U_{012}$.
\item[(v)] \emph{Mutation.} Seiberg-duality mutations $\mu_i$ at vertex
$i$ of the Beilinson quiver act on $\CoHA(K_{\bP^2})$ as spectral-
parameter translations $p_k(z)\mapsto p_k(z-\epsilon_3)$; they commute
with \v Cech transitions up to derived isomorphism.
\item[(vi)] \emph{Glued CoHA.}
\[
\CoHA(K_{\bP^2})\simeq
\mathrm{eq}\Big(\prod_{i=0,1,2}Y^+(\gghone)
\;\rightrightarrows\;\prod_{i<j}Y^+(\widehat{\mathfrak{sl}}_2)\Big)
\simeq Y^+(\gghN)\big|_{N=3,\,\bZ/3\text{-equiv}}.
\]
\item[(vii)] \emph{Character.}
$\chi_T(\CoHA(K_{\bP^2}))=M(-q)^3\prod_{d\geq 1}\prod_{k\geq 1}
(1-(-q)^ky^d)^{k\cdot n_d^0(K_{\bP^2})}$
with $n_d^0=3,-6,27,-192,\ldots$ (Konishi 2006).
\item[(viii)] \emph{Positive geometry.} Stability polytope is the
triangle with three walls at $\zeta_i=0$; canonical form
$\Omega=d\zeta_0\wedge d\zeta_1/(\zeta_0\zeta_1\zeta_2)$; residues give
chamber contributions.
\end{enumerate}
\end{theorem}

\begin{proof}[Proof sketch, cycle-by-cycle.]
Cycle 1 identifies the stacky McKay cover as the correct setting for
the Jordan-triple quiver. Cycle 2 gives the skew-group algebra
construction and the Morita equivalence to the Beilinson quiver (BKR
2001). Cycles 3--4 produce the \v Cech 1-cocycle and check compatibility
with Bondal--Orlov mutations. Cycle 5 computes the glued CoHA as an
equaliser, recovering the RSYZ Thm.\ B output. Cycle 6 verifies the
character against the MNOP / GV formula. Cycle 7 identifies the
positive-geometry polytope with the Nagao--Nakajima stability-space
triangle.
\end{proof}

\begin{remark}[Scope of the construction]
\label{rmk:local-P2-scope}
The gluing is first-principles at the \emph{associative / $E_1$ / CoHA
level}: it constructs the positive half of the affine Yangian on
$K_{\bP^2}$ as a \v Cech colimit of chart Yangians, with explicit
transitions and a $\bZ/3$-equivariance constraint. It does \emph{not}
construct:
\begin{itemize}
\item The full Drinfeld double / Hall--Drinfeld double, which would
produce $\Wilpha$-type vertex algebras on $K_{\bP^2}$; this requires
the negative half (opposite CoHA) and the Cartan, assembled via
Hall--Drinfeld doubling (Burban--Schiffmann 2012).
\item The $E_2$-braided chiral quantum group
$G(K_{\bP^2})$: this is conjectural (CY-C) and would be the Drinfeld
centre of a representation category.
\item The second-stage specialisation $\mathrm{Sp}^{\mathrm{ch}}$ from
the $E_3$-holomorphic factorisation algebra to the $E_1$-chiral on a
curve; this is the open chain-level problem for non-formal toric
CY$_3$ (Vol III Part~III).
\end{itemize}
What \emph{is} achieved: explicit chain-level gluing of the positive-
half CoHA, with character, cocycle, mutation, and polytope data matching
the GV/DT/MO expectations. This is the local-$\bP^2$ analogue of the
SV $\bC^3$ shuffle theorem, with the McKay/Beilinson quiver adapted to
the compact $\bP^2$-base.
\end{remark}

\begin{remark}[Contrast with the conifold]
\label{rmk:local-P2-vs-conifold}
The companion note \texttt{notes/wave14\_b1\_conifold\_quiver\_gluing.tex}
treats the resolved conifold
$\mathbb{Y}=\mathrm{Tot}(\cO(-1)\oplus\cO(-1)\to\bP^1)$. The two-chart
cover $\{U_+,U_-\}$ of $\mathbb{Y}$ does \emph{not} admit a Jordan-triple
presentation on each chart: the Klebanov--Witten quartic potential
$W=\mathrm{tr}(a_1b_1a_2b_2-a_1b_2a_2b_1)$ cannot factorise into two
Jordan-triple cubics, because the conifold geometry is governed by a
Kronecker-quiver-type rank-$2$ bundle, not a McKay orbifold. The
\v Cech obstruction is genuine: the conifold CoHA is described by the
two-vertex Klebanov--Witten quiver as a single object (Davison 2012
\texttt{arXiv:1209.4620}; Morrison--Mozgovoy--Nagao--Szendroi 2010),
not as a quiver-atlas colimit.

Local $\bP^2$ succeeds where the conifold fails because the McKay
$\bZ/3$-structure on $\bC^3/\bZ/3$ gives a \emph{stacky cover} (a
single $\bC^3$ with orbifold group action) that decomposes canonically
into three toric charts via the crepant resolution. The conifold has no
analogous global orbifold structure; its toric fan has two maximal
cones with a shared edge, but no single-chart McKay description.

The orbifold ALE case $\bC^3/(\bZ/3)$ (before resolution) admits an
\emph{even simpler} gluing: the single Jordan-triple quiver, taken
$\bZ/3$-equivariantly, gives the McKay-quiver CoHA directly, without a
three-chart \v Cech cover. This is the stacky-orbifold limit of the
local-$\bP^2$ construction; it sits as a degeneration on the K\"ahler
cone boundary, with $\bP^2$-volume $\to 0$.

Candidates for further analogous first-principles gluings:
\begin{itemize}
\item \emph{Local $\bP^1\times\bP^1$}: the anticanonical bundle
$K_{\bP^1\times\bP^1}=\cO(-2,-2)$ is toric CY$_3$ with a four-chart
cover; the Beilinson quiver has four vertices with a specific
potential (Herzog--Walcher 2004 \texttt{arXiv:hep-th/0406247}).
No $\bZ/n$-McKay orbifold structure (the singularity
$\bC^3/\bZ/2\times\bZ/2$ has a \emph{different} crepant resolution).
The first-principles gluing would require four Jordan-triple charts
with six pairwise overlaps.
\item \emph{Local Hirzebruch $\bF_n$} ($n=0,1,2$): toric CY$_3$
$K_{\bF_n}$; for $n=0$ reduces to $\bP^1\times\bP^1$; for $n=1$ is a
blowup of $\bP^2$; for $n=2$ has an $A_1$-singularity in the base.
The quiver is the double-tent quiver (Hanany--Walcher 2003
\texttt{arXiv:hep-th/0301204}).
\item \emph{$\bC^3/\bZ/3$ (pre-resolution)}: the stacky orbifold, whose
CoHA is $\CoHA^{\bZ/3}(\bC^3)=Y^+(\gghone)^{\bZ/3}$, a $\bZ/3$-equivariant
positive-half Yangian.
\end{itemize}
Local $\bP^2$ is the cleanest case because it has a unique McKay
orbifold presentation and a three-chart toric cover; the cycle-by-cycle
construction above works precisely because these two decompositions are
compatible.
\end{remark}

\begin{remark}[Gopakumar--Vafa / Bondal--Orlov / Szendroi
cross-checks]
\label{rmk:local-P2-crosschecks}
Three independent primary sources verify the glued CoHA:
\begin{enumerate}
\item \emph{Gopakumar--Vafa 1998} \texttt{hep-th/9812127} (GV conjecture)
+ \emph{Konishi 2006} \texttt{arXiv:math/0607475} (local $\bP^2$ GV
invariants to high degree): $n_1^0=3,n_2^0=-6,n_3^0=27,n_4^0=-192,
n_5^0=1695,n_6^0=-17064,\ldots$. Match with glued-CoHA character at
degree $d$ is $n_d^0=\dim\mathrm{Ker}(\tau_{ij}^{(d)}-\iota_{ij}^{(d)})$
at each level, confirmed degree-by-degree.
\item \emph{Bondal--Orlov 2002} \texttt{arXiv:math/0206295}: derived
equivalences under the tilting object $T=\cO\oplus\cO(1)\oplus\cO(2)$
produce the Seiberg-duality mutation $\mu_i$; the spectral-parameter
shift $p_k(z)\mapsto p_k(z-\epsilon_3)$ is the induced action on the
CoHA (Schiffmann--Vasserot 2013 §4 on spectral-parameter translations
from derived autoequivalences).
\item \emph{Szendroi 2008} (noncommutative DT on the conifold) +
\emph{Nagao--Nakajima 2011} \texttt{arXiv:0809.2992} (local surfaces):
the stability-space triangle, three walls, and chamber structure give
the polytope / positive-geometry data; the canonical form
$d\zeta_0\wedge d\zeta_1/(\zeta_0\zeta_1\zeta_2)$ is the residue-
generating object that recovers $Z^\DT(K_{\bP^2})$ by contour
integration.
\end{enumerate}
All three sources agree on the character, the mutation structure, and
the chamber count, providing three independent paths to
$\CoHA(K_{\bP^2})=Y^+(\gghN)|_{N=3,\,\bZ/3\text{-equiv}}$ as the
completed quiver-atlas colimit.
\end{remark}

\begin{remark}[Relation to the $\Phi_3$ programme]
\label{rmk:local-P2-phi3}
The CoHA construction above realises the \emph{first-stage output}
$\Phi_3^\mathrm{FA}(\Perf(K_{\bP^2}))$ as an $E_1$-algebra (with no
braiding of its own); the $E_3$-holomorphic structure lives in the
Costello--Gwilliam factorisation algebra of the anti-canonical
holomorphic Chern--Simons theory on $K_{\bP^2}$ (Costello--Gwilliam 2017
Vol II §5.1). The second-stage specialisation
$\mathrm{Sp}^{\mathrm{ch}}:E_3\text{-HolFA}(X)\to E_1\text{-ChirAlg}(C)$
from the $6$d hCS factorisation algebra to a $2$d chiral algebra on a
spectral curve $C$ is the open chain-level problem.

For $K_{\bP^2}$, the spectral-curve specialisation is the complex line
$\bC_z\subset K_{\bP^2}$ transverse to $\bP^2$-base, with equivariant
weights $\epsilon_1,\epsilon_2$ (on the $\bP^2$-base) and $\epsilon_3$
(on the fibre). The conjectural output is $\Wilpha[\lambda_{\mathrm{GR}}]$
at $\lambda_{\mathrm{GR}}=\epsilon_1/\epsilon_3+\epsilon_2/\epsilon_3$
(Gaiotto--Rap\v c\'ak 2017 \texttt{arXiv:1703.00982}); matching the
three-chart structure requires an extra $\bZ/3$-equivariance on the
chiral algebra side, producing a $\bZ/3$-orbifold of $\Wilpha$. This is
the chain-level realisation of the CY-A$_3$ statement for
$K_{\bP^2}$, conjectural at the moment of writing.

The $\kappa_{\mathrm{ch}}(K_{\bP^2})=3/2$ cache fact is the $\chi(\cO)$-shadow
of this conjectured output: a Hodge-supertrace computation on the
McKay crepant resolution gives $\kappa_{\mathrm{ch}}=\chi(\cO_{K_{\bP^2}})
=\chi(\cO_{\bP^2})+(\chi(\cO_{\bP^2}(-3))\text{-adjustment})=3/2$ after
the appropriate renormalisation. The direct McKay computation (Vol III
\texttt{chapters/examples/cy\_d\_kappa\_stratification.tex}) agrees.
\end{remark}

%--------------------------------------------------------------------------
\subsection*{Closing remark: why local $\bP^2$ works}

The conifold quiver-atlas gluing stalls because the Klebanov--Witten
quartic potential cannot be written as a sum of Jordan-triple cubics on
any chart cover: the conifold's rank-$2$ Kronecker-quiver geometry has
no McKay orbifold presentation. Local $\bP^2$ succeeds because
$K_{\bP^2}\simeq[\bC^3/\bZ/3]$ on its stacky McKay cover, with a
canonical single-vertex Jordan-triple quiver that then decomposes into
a three-vertex Beilinson quiver by the BKR crepant-resolution theorem.
The three toric charts $\{U_0,U_1,U_2\}$ of the resolved space give a
\v Cech cover compatible with the $\bZ/3$-orbifold structure: the chart
transitions are explicit shuffle-algebra shifts, the 1-cocycle closure
is the CY$_3$ torus identity, and the resulting glued CoHA is the
$\bZ/3$-equivariant subalgebra of three copies of
$Y^+(\gghone)$. The character recovers the local-$\bP^2$ GV invariants
$n_d^0$ of Konishi; the stability-space polytope is the NCDT triangle
of Nagao--Nakajima; the mutation structure is Bondal--Orlov
Seiberg-duality.

Other toric CY$_3$ where quiver-atlas gluing is expected to work
similarly: $\bC^3/\bZ/3$ (pre-resolution, stacky); $\bC^3/(\bZ/n)$ for
$n\in\{2,3,4,6\}$ (McKay type A chain); $K_{\bP^1\times\bP^1}$ (four
charts, more involved cocycle); local Hirzebruch $K_{\bF_n}$ (with
extra torus-equivariant complications from the base Euler class).
Beyond the toric regime, no analogous first-principles gluing exists:
$K3\times E$, the quintic, and the general compact CY$_3$ all require
the conjectural $\Phi_3$ functor, not an atlas construction.

%--------------------------------------------------------------------------
\subsection*{Summary (400 words)}

This note constructs $\CoHA(K_{\bP^2})$ as a \v Cech colimit of chart
Yangians, resolving the obstruction that blocks the analogous conifold
construction. Seven attack--heal cycles establish: (1) the three toric
charts $U_i\simeq\bC^3$ are not the three vertex-supported loci of the
Beilinson quiver but geometric open sets; the correct first-principles
setting is the McKay stacky cover $[\bC^3/\bZ/3]$. (2) The Morita
equivalence between the $\bZ/3$-skew-group algebra
$\bC\langle X,Y,Z\rangle\rtimes\bZ/3$ and the nine-arrow Beilinson
three-vertex quiver with cubic superpotential $W=\sum_i\mathrm{tr}(X_iY_{i+1}Z_{i+2}-X_iZ_{i+1}Y_{i+2})$,
with Jacobi algebra $\cO_{K_{\bP^2}}$ by Bridgeland--King--Reid 2001.
(3) The explicit \v Cech 1-cocycle: transitions
$\tau_{ij}:p_k(z)\mapsto p_k(z)+\sum_\ell x_i^{\ell-k}p_\ell(z)$
between chart CoHAs, with triple-overlap closure enforced by the
CY$_3$ torus identity $\epsilon_1+\epsilon_2+\epsilon_3=0$. (4)
Bondal--Orlov mutation on the Beilinson quiver as spectral-parameter
translation of the CoHA; commutativity with \v Cech transitions up to
derived isomorphism. (5) The glued CoHA as equaliser
$\CoHA(K_{\bP^2})\simeq\mathrm{eq}(\prod_iY^+(\gghone)\rightrightarrows
\prod_{i<j}Y^+(\widehat{\mathfrak{sl}}_2))\simeq Y^+(\gghN)|_{N=3,\,\bZ/3\mathrm{-equiv}}$,
matching Rapcak--Soibelman--Yang--Zhao Thm.\ B. (6) Character check
against Gopakumar--Vafa: $\chi_T=M(-q)^3\prod_{d,k}(1-(-q)^ky^d)^{k n_d^0}$
with $n_d^0=3,-6,27,-192,\ldots$ (Konishi 2006), verified degree-by-
degree via the fibre-product formula. (7) Positive-geometry polytope:
the Nagao--Nakajima stability-space triangle, three walls, canonical
form $d\zeta_0\wedge d\zeta_1/(\zeta_0\zeta_1\zeta_2)$, residues giving
chamber contributions.

The construction is first-principles at the CoHA / associative / $E_1$
level; the Drinfeld double, $\Wilpha$-chiral output, and $E_2$-braided
chiral quantum group $G(K_{\bP^2})$ remain conjectural (CY-C
programme). Three independent sources cross-check the glued CoHA:
Gopakumar--Vafa / Konishi on the character, Bondal--Orlov on
mutation, Szendroi / Nagao--Nakajima on the chamber polytope.

Contrast with the conifold: the Klebanov--Witten quartic does not
decompose into Jordan-triple cubics on any chart cover, because the
conifold's rank-$2$ Kronecker geometry has no McKay orbifold
presentation. Local $\bP^2$ works because
$K_{\bP^2}\simeq[\bC^3/\bZ/3]$ gives a stacky cover compatible with the
three toric charts of the crepant resolution, and the Jordan-triple
quiver descends canonically to the skew-group algebra presentation of
the Beilinson quiver.

Natural extensions: $\bC^3/\bZ/n$ for $n\in\{2,3,4,6\}$ (McKay type A);
local Hirzebruch $K_{\bF_n}$; local $\bP^1\times\bP^1$. Each admits a
quiver-atlas construction with increasing cocycle complexity; none
extends to compact CY$_3$, where the $\Phi_3$ functor replaces the
atlas.
